% ID: MAT-TRI-EQG-013
% Titolo: Valutare affermazioni goniometriche (1)
% Disciplina: Matematica
% Area: trigonometria
% Argomento: equazioni_goniometriche
% Classe: Quarta secondaria di secondo grado
% Difficolta: 3
% Tipo: vero_falso
% Tag: matematica, trigonometria, equazioni_goniometriche
% Fonte: verifica_goniometria_base_anonimizzata
% Provenienza: materiale fornito dal docente; testo anonimizzato; soluzione redatta per Exergo

\begin{esercizio}
Vero o falso? 

\begin{enumerate}
\item L'equazione $\sin x=\dfrac{4}{3}$ è impossibile. 

\item L'equazione $\tan x=0$ ammette una sola soluzione
nell'intervallo $[0;2\pi)$. 

\item Il valore $x=\dfrac{3\pi}{4}$ è soluzione dell'equazione
$\dfrac{\sqrt{2}\sin x-1}{\cos x+\dfrac{\sqrt{2}}{2}}=0$.


\item L'equazione $\cos 2x=1$ ammette esattamente due soluzioni
nell'intervallo $[0;2\pi)$. 
\end{enumerate}
Salvo diversa indicazione, lavora in $\mathbb R$; angoli in radianti.
\end{esercizio}

\begin{soluzione}
a) Vero: il seno appartiene a $[-1,1]$.
b) Falso: le soluzioni sono $0$ e $\pi$.
c) Falso: il denominatore si annulla per $x=3\pi/4$.
d) Vero: $\cos2x=1$ equivale a $x=k\pi$, quindi $0$ e $\pi$.
\end{soluzione}
